\chapter{Banco ampliado de test teóricos (edición de repaso)}
\label{ch:test-matmo-2026}
\noindent\textbf{Nota editorial.} Preguntas originales para practicar conceptos del temario de la ULE 2026--2027. Inspiradas por los temas de las fuentes de estudio aportadas, pero no reproducen sus enunciados. Las respuestas de apuntes de estudiantes pueden ser erróneas. No son preguntas oficiales, no certifican vigencia normativa ni sustituyen las fuentes citadas en el manual.

\noindent\textbf{Uso:} 36 preguntas, una opción correcta por pregunta. Conteste primero sin consultar el solucionario. Justifique además por qué descarta las alternativas.

\section*{Materiales, microestructura y aleaciones}
\subsection*{T01. ¿Cuál es la familia principal de la aleación 8090?}
\begin{description}
\item[a)] Aleación base níquel
\item[b)] Aleación aluminio-litio
\item[c)] Acero maraging
\end{description}

\subsection*{T02. ¿Qué implica que un componente trabaje por debajo de 200 grados Celsius?}
\begin{description}
\item[a)] Que deben excluirse todos los metales
\item[b)] Que debe usarse obligatoriamente una resina termoestable
\item[c)] Que la temperatura por sí sola no selecciona el material
\end{description}

\subsection*{T03. ¿Qué permite el envejecimiento artificial de aleaciones de aluminio tratables?}
\begin{description}
\item[a)] Controlar la precipitación y modificar sus propiedades
\item[b)] Convertir aluminio en titanio
\item[c)] Aumentar solamente la dureza superficial por difusión de carbono
\end{description}

\subsection*{T04. En una lámina reforzada con fibras, ¿qué función cumple la matriz?}
\begin{description}
\item[a)] Nunca soporta tensión
\item[b)] Transfiere carga entre fibras y contribuye a protección y respuesta transversal
\item[c)] Reemplaza el refuerzo cuando hay cargas longitudinales
\end{description}

\subsection*{T05. ¿Cuál es la distinción general entre termoplástico y termoestable?}
\begin{description}
\item[a)] Ambos se vuelven líquidos reversiblemente al calentar
\item[b)] Todo termoestable es un termoplástico muy blando
\item[c)] Los termoestables curados tienen redes reticuladas y no se funden como termoplásticos
\end{description}

\subsection*{T06. ¿Qué estructuras cristalinas corresponden a las fases principales del titanio?}
\begin{description}
\item[a)] Alfa HCP y beta BCC
\item[b)] Alfa BCC y beta FCC
\item[c)] Alfa FCC y beta HCP
\end{description}

\subsection*{T07. ¿Qué es GLARE?}
\begin{description}
\item[a)] Un acero inoxidable reforzado con carbono
\item[b)] Un laminado híbrido con capas de aluminio y material compuesto de fibra de vidrio
\item[c)] Un polímero elastomérico sin metal
\end{description}

\subsection*{T08. Si un adhesivo falla dentro de su propio espesor, el fallo es:}
\begin{description}
\item[a)] Adhesivo interfacial
\item[b)] Del sustrato
\item[c)] Cohesivo del adhesivo
\end{description}

\subsection*{T09. En un panel sándwich sometido a flexión, la función habitual del núcleo es:}
\begin{description}
\item[a)] Separar las caras, transmitir cortante y ayudar frente a inestabilidad
\item[b)] Proporcionar por sí solo toda la rigidez axial
\item[c)] Evitar toda deformación de las caras
\end{description}

\subsection*{T10. ¿Para qué puede utilizarse una superaleación monocristalina en un álabe?}
\begin{description}
\item[a)] Para evitar toda necesidad de inspección
\item[b)] Para reducir mecanismos asociados a límites de grano a alta temperatura
\item[c)] Para garantizar que el punto de fusión supera 2000 grados Celsius
\end{description}

\section*{Fractura, fatiga, fluencia y ensayos}
\subsection*{T11. Para una barra a tracción dimensionada por resistencia, ¿qué índice de material favorece menor masa bajo hipótesis constantes?}
\begin{description}
\item[a)] Densidad alta
\item[b)] Solo dureza superficial
\item[c)] Resistencia admisible dividida por densidad
\end{description}

\subsection*{T12. En mecánica de fractura lineal elástica, ¿qué comparación define el criterio idealizado de fractura?}
\begin{description}
\item[a)] Factor de intensidad aplicado frente a tenacidad aplicable
\item[b)] Módulo de Young frente a densidad
\item[c)] Dureza frente a temperatura de fusión
\end{description}

\subsection*{T13. En la ecuación de una grieta central pasante de longitud total 2a, ¿qué representa a?}
\begin{description}
\item[a)] El espesor de la placa
\item[b)] La semilongitud de la grieta
\item[c)] El radio de curvatura de la punta
\end{description}

\subsection*{T14. En la ley de Paris, la fuerza impulsora cíclica se expresa con:}
\begin{description}
\item[a)] La tensión media exclusivamente
\item[b)] El límite elástico del material
\item[c)] El rango de intensidad Delta K
\end{description}

\subsection*{T15. ¿Qué limitación básica tiene la regla lineal de Miner?}
\begin{description}
\item[a)] No representa por sí misma secuencia ni interacción entre bloques
\item[b)] Incluye siempre efectos ambientales y de superficie
\item[c)] Garantiza vida segura cuando D es menor que uno
\end{description}

\subsection*{T16. Sobre fluencia en un monocristal metálico, ¿qué afirmación es válida?}
\begin{description}
\item[a)] Solo puede haber fluencia por difusión
\item[b)] También pueden intervenir mecanismos de dislocaciones y otros
\item[c)] La fluencia es imposible sin bordes de grano
\end{description}

\subsection*{T17. El parámetro de Larson-Miller se utiliza para:}
\begin{description}
\item[a)] Eliminar ensayos de rotura
\item[b)] Determinar sin hipótesis la vida de cualquier material
\item[c)] Correlacionar tiempo de rotura y temperatura en un dominio de calibración
\end{description}

\subsection*{T18. Para evaluar una discontinuidad interna por ultrasonidos, debe considerarse:}
\begin{description}
\item[a)] Velocidad de onda, calibración y geometría del eco
\item[b)] Únicamente el color de la superficie
\item[c)] Que cualquier eco revela directamente el tamaño verdadero
\end{description}

\subsection*{T19. La termografía infrarroja detecta principalmente:}
\begin{description}
\item[a)] Solamente transmisión de luz visible
\item[b)] Contrastes de radiación térmica relacionados con temperatura y emisividad
\item[c)] La composición exacta de una aleación sin calibración
\end{description}

\subsection*{T20. En corrosión galvánica, ¿qué factor es relevante?}
\begin{description}
\item[a)] Que ambos materiales tengan el mismo nombre comercial
\item[b)] Únicamente el espesor del recubrimiento
\item[c)] Contacto eléctrico, electrolito y potenciales relativos
\end{description}

\section*{Normativa y mantenimiento}
\subsection*{T21. En la Unión Europea, una AD aplicable obliga a:}
\begin{description}
\item[a)] Cumplir lo requerido en su ámbito y plazos
\item[b)] Aceptar cualquier revisión posterior de un SB como equivalente
\item[c)] Ignorar configuración y número de serie
\end{description}

\subsection*{T22. La aprobación Parte 145 se refiere principalmente a:}
\begin{description}
\item[a)] El certificado de tipo de motor
\item[b)] La organización de mantenimiento y su alcance aprobado
\item[c)] La licencia de piloto
\end{description}

\subsection*{T23. La Parte 66 regula principalmente:}
\begin{description}
\item[a)] El diseño de piezas
\item[b)] El programa de fabricación
\item[c)] Licencias de mantenimiento y privilegios asociados
\end{description}

\subsection*{T24. La Parte 147 aborda principalmente:}
\begin{description}
\item[a)] Organizaciones aprobadas de formación y exámenes de mantenimiento
\item[b)] El permiso de vuelo de una aeronave dañada
\item[c)] El seguimiento de corrosión mediante ultrasonidos
\end{description}

\subsection*{T25. ¿Qué función caracteriza a una organización CAMO?}
\begin{description}
\item[a)] Expedir siempre certificados de tipo
\item[b)] Gestionar la aeronavegabilidad continuada dentro de sus privilegios
\item[c)] Sustituir automáticamente el trabajo de un taller Parte 145
\end{description}

\subsection*{T26. ¿Qué documento se asocia a una revisión de aeronavegabilidad?}
\begin{description}
\item[a)] IPC
\item[b)] CMM
\item[c)] ARC
\end{description}

\subsection*{T27. En normativa europea, Parte 21 se vincula sobre todo a:}
\begin{description}
\item[a)] Certificación de diseño y producción y procedimientos de aeronavegabilidad inicial
\item[b)] Cálculo de resistencia a fatiga por Miner
\item[c)] Clasificación de grietas por partículas magnéticas
\end{description}

\subsection*{T28. ¿Qué par de reglamentos se relaciona con aeronavegabilidad inicial y continuada, respectivamente?}
\begin{description}
\item[a)] 1321/2014 y 748/2012
\item[b)] 748/2012 y 1321/2014
\item[c)] 2018/1139 y una AD aislada
\end{description}

\subsection*{T29. Un certificado de tipo acredita principalmente:}
\begin{description}
\item[a)] Que todo ejemplar individual puede volar sin inspecciones
\item[b)] Que una inspección de taller ha concluido
\item[c)] Un diseño de tipo aprobado conforme a su base de certificación
\end{description}

\subsection*{T30. Para aplicar Parte ML a un avión, hay que comprobar:}
\begin{description}
\item[a)] Categoría, masa, complejidad, operación y excepciones del ámbito
\item[b)] Solo que sea monomotor
\item[c)] Solo que tenga menos de veinte años
\end{description}

\subsection*{T31. Si una tarea del AMP vence por horas o calendario, lo que ocurra antes:}
\begin{description}
\item[a)] Se usa siempre la fecha que más tarde llegue
\item[b)] Debe respetarse el primer límite aplicable
\item[c)] Es una recomendación sin control documental
\end{description}

\subsection*{T32. ¿Cómo se interpreta una tabla de AD de FAA para una aeronave bajo ámbito EASA?}
\begin{description}
\item[a)] Se supone automáticamente aplicable sin comprobar matrícula y régimen
\item[b)] Se descarta siempre por ser americana
\item[c)] Se estudian autoridad, aplicabilidad, adopción y estado de cada AD
\end{description}

\subsection*{T33. Si un PDF BOE se titula como si fuese norma de aeronavegabilidad, antes de citarlo debe:}
\begin{description}
\item[a)] Comprobar el identificador, fecha, objeto y disposiciones efectivas
\item[b)] Confiar únicamente en el nombre del archivo
\item[c)] Tratar cualquier BOE como un Reglamento de aeronavegabilidad
\end{description}

\subsection*{T34. Cuando se actualiza una norma, para contestar un examen hay que distinguir:}
\begin{description}
\item[a)] Solo el día de descarga de la web
\item[b)] Fecha de publicación, entrada en vigor, aplicación y transitorios
\item[c)] Únicamente el número de páginas
\end{description}

\subsection*{T35. Si un técnico dispone de licencia B1.1 pero carece de evidencia de habilitación y autorización exigidas:}
\begin{description}
\item[a)] Su licencia basta para todo trabajo de ese tipo
\item[b)] Puede certificar cualquier trabajo por experiencia verbal
\item[c)] No se han demostrado todos los privilegios necesarios
\end{description}

\subsection*{T36. Una tendencia de EGT anómala en seguimiento de motor:}
\begin{description}
\item[a)] Debe analizarse con condición operativa, calidad de datos y otros indicadores
\item[b)] Determina por sí sola qué pieza debe retirarse
\item[c)] Permite extender automáticamente vida límite
\end{description}

